\documentclass[conference,twocolumn,spanish]{IEEEtran}

\usepackage[utf8]{inputenc}
\usepackage[T1]{fontenc}
\usepackage[spanish,es-tabla]{babel}
\usepackage{amsmath}
\usepackage{graphicx}
\usepackage{fancyhdr}
\usepackage{booktabs}
\usepackage{tabularx}
\usepackage{pgfplots}
\pgfplotsset{compat=1.18}
\usepackage{xurl}
\usepackage{hyperref}
\usepackage{cite}

% Metadatos internos del documento PDF (incluye derechos y licencia CC BY-NC-ND 4.0)
\hypersetup{
    pdftitle={Boletín | Centro de Inteligencia Artificial e Innovación -- 11 al 18 de noviembre de 2026},
    pdfauthor={Miguel Ángel Osorio Orduz; Jhoan Sebastian Torres Garcia; Frank Camilo Zuñiga Vargas; Docentes: Carlos Alfonso Pineda Ortiz, Diego Alexander Cruz Melo},
    pdfsubject={Boletín Semanal de Vigilancia Tecnológica en Inteligencia Artificial (CIAI - UTS)},
    pdfkeywords={Inteligencia Artificial, LLM, CIAI, UTS, Vigilancia Tecnológica, Benchmarks, GPT-6.2 Astra, GPT-6.2 Sol, Qwen4 Max, Gemini 4 Hydrogen, Kitesurf, NVIDIA, STAR 50},
    pdfcreator={LaTeX con clase IEEEtran},
    pdfproducer={Centro de Inteligencia Artificial e Innovación (CIAI) - Unidades Tecnológicas de Santander},
    pdfcopyright={Copyright (C) 2026 CIAI - UTS. Licencia Creative Commons Atribución-NoComercial-SinDerivadas 4.0 Internacional (CC BY-NC-ND 4.0).},
    pdflicenseurl={https://creativecommons.org/licenses/by-nc-nd/4.0/deed.es}
}

% Configuración de encabezado y pie de página
% Nota: si compila localmente dentro de ediciones/2026-11-semana2/, ajusta la ruta del logo a ../../assets/uts_logonegro.png
\pagestyle{fancy}
\fancyhf{}
\fancyhead[L]{\small Boletín Semanal CIAI -- 11 al 18 de noviembre de 2026}
\fancyhead[R]{\IfFileExists{assets/uts_logonegro.png}{\includegraphics[height=0.8cm]{assets/uts_logonegro.png}}{\textbf{CIAI -- UTS}}}
\fancyfoot[L]{\scriptsize \copyright~2026 CIAI -- UTS | Licencia \href{https://creativecommons.org/licenses/by-nc-nd/4.0/deed.es}{CC BY-NC-ND 4.0} (Uso académico, no comercial, sin derivadas)}
\fancyfoot[R]{\scriptsize \thepage}
\renewcommand{\headrulewidth}{0.4pt}
\renewcommand{\footrulewidth}{0.2pt}

% Macro dinámica para gráficas de avance (0% a 100%) de los Proyectos CIAI
% Uso: \graficaavanceproyecto{Título corto}{Hito4,Hito3,Hito2,Hito1}{{(Hito4,0) (Hito3,25) (Hito2,70) (Hito1,100)}}
\newcommand{\graficaavanceproyecto}[3]{%
\begin{figure}[htbp]
\centering
\begin{tikzpicture}
\begin{axis}[
    xbar,
    width=0.92\columnwidth,
    height=3.6cm,
    bar width=7pt,
    xlabel={Porcentaje de avance (\%)},
    xmin=0,
    xmax=114,
    xtick={0,25,50,75,100},
    symbolic y coords={#2},
    ytick=data,
    tick label style={font=\scriptsize},
    label style={font=\scriptsize},
    xmajorgrids=true,
    grid style={dashed, gray!35},
    nodes near coords={\pgfmathprintnumber{\pgfplotspointmeta}\%},
    every node near coord/.append style={font=\tiny, anchor=west}
]
\addplot[fill=green!55!black, draw=green!70!black] coordinates {#3};
\end{axis}
\end{tikzpicture}
\caption{Porcentaje de avance por hitos: #1.}
\end{figure}
}

\title{Boletín | Centro de Inteligencia Artificial e Innovación\\
       \large Edición Semanal: 11 al 18 de noviembre de 2026}

\author{
  \IEEEauthorblockN{Miguel Ángel Osorio Orduz}
  \IEEEauthorblockN{Jhoan Sebastian Torres Garcia}
  \IEEEauthorblockN{Frank Camilo Zuñiga Vargas}
  \IEEEauthorblockA{Docentes: Carlos Alfonso Pineda Ortiz, Diego Alexander Cruz Melo}
  \IEEEauthorblockA{
    Centro de Inteligencia Artificial e Innovación (CIAI)\\
    Unidades Tecnológicas de Santander --- Colombia\\
    mangelosorio@uts.edu.co\\
    jhoanstorres@uts.edu.co\\
    fczuniga@uts.edu.co
  }
}

\begin{document}

\maketitle
\thispagestyle{fancy}

\begin{abstract}
Esta edición del Boletín del Centro de Inteligencia Artificial e Innovación (CIAI) sintetiza los avances clave en inteligencia artificial entre el 11 y el 18 de noviembre de 2026. En investigación se analizan las incorporaciones al \textit{Artificial Analysis Intelligence Index v5.0.1} tras el lanzamiento de GPT-6.2~Astra, GPT-6.2~Sol, Qwen4~Max y Gemini~4~Hydrogen, acompañadas de comparativas verticales de rendimiento, tarifas API (\textit{Input/Output}) y latencia inicial (\textit{TTFT}). En herramientas se documentan el motor de navegación agéntica Kitesurf de Cloudflare y la migración de \textit{Gems} a \textit{Skills} en Google Gemini. En noticias se cubren la demanda de California contra OpenAI, la recta final de la oferta pública inicial de Anthropic en Nasdaq, los nuevos controles estadounidenses sobre el acceso remoto a semiconductores y la evolución bursátil de NVIDIA y del índice tecnológico chino SSE~STAR~50.
\end{abstract}

\begin{IEEEkeywords}
Inteligencia Artificial, LLM, CIAI, UTS, Vigilancia Tecnológica, Artificial Analysis, GPT-6.2 Astra, GPT-6.2 Sol, Qwen4 Max, Gemini 4 Hydrogen, Kitesurf, NVIDIA, SSE STAR 50
\end{IEEEkeywords}

% =========================================================
\section{Investigación}
% =========================================================

Durante la semana del 11 al 18 de noviembre de 2026, el ranking \textit{Artificial Analysis Intelligence Index v5.0.1} registró una renovación significativa en la frontera de desempeño y eficiencia operativa~\cite{aa_index_2026}.

El 17 de noviembre, OpenAI presentó los modelos \textbf{GPT-6.2~Astra} y \textbf{GPT-6.2~Sol}, ambos con ventana de contexto de $1\,\text{M}$ de tokens~\cite{openai_gpt62_2026}. En esfuerzo de razonamiento máximo (\textit{Max}), \texttt{GPT-6.2~Astra} debutó en el segundo puesto de la clasificación general con $62$ puntos, un costo promedio de $\text{USD}~4{,}00$ por tarea ($\text{USD}~0{,}66$ en esfuerzo bajo), tarifas API de $\text{USD}~8{,}00$ y $\text{USD}~40{,}00$ por millón de tokens de entrada y salida, un tiempo medio de $12{,}9\,\text{min}$ por tarea y una latencia al primer token de respuesta (\textit{TTFT}) de $147{,}65\,\text{s}$~\cite{aa_index_2026,openai_gpt62_2026}. Por su parte, \texttt{GPT-6.2~Sol (Max)} alcanzó $57$ puntos con un costo de $\text{USD}~1{,}01$ por tarea ($\text{USD}~0{,}18$ en esfuerzo bajo), tarifas de $\text{USD}~2{,}00$ / $\text{USD}~10{,}00$, $8{,}0\,\text{min}$ por tarea y un \textit{TTFT} de $66{,}70\,\text{s}$, posicionándose como referente de eficiencia económica~\cite{aa_index_2026,openai_gpt62_2026}.

El 16 de noviembre, Alibaba liberó \textbf{Qwen4~Max}, que obtuvo $43$ puntos en el índice con una ventana de contexto de $984\,\text{k}$ tokens, costo promedio de $\text{USD}~3{,}27$ por tarea, tarifas API de $\text{USD}~1{,}20$ / $\text{USD}~6{,}00$, velocidad de generación de $54\,\text{tok/s}$ y un \textit{TTFT} de $38{,}40\,\text{s}$~\cite{aa_index_2026,alibaba_qwen4_2026}. Previamente, el 12 de noviembre, Google incorporó \textbf{Gemini~4~Hydrogen}, alcanzando $54$ puntos con contexto de $1\,\text{M}$ de tokens, tarifas API de $\text{USD}~1{,}50$ / $\text{USD}~7{,}50$ ($\text{USD}~3{,}16$ por tarea en esfuerzo alto), velocidad de $173\,\text{tok/s}$ y un \textit{TTFT} de $44{,}89\,\text{s}$~\cite{aa_index_2026,google_gemini4_hydrogen_2026}.

\begin{table}[htbp]
\centering
\caption{Comparativa de modelos líderes en Artificial Analysis Intelligence Index v5.0.1 (al 18 de noviembre de 2026)~\cite{aa_index_2026}.}
\label{tab:aa_modelos_sem2}
\scriptsize
\begin{tabularx}{\columnwidth}{@{}l l c c c@{}}
\toprule
\textbf{Modelo} & \textbf{Proveedor} & \textbf{Índice} & \textbf{Costo/tarea} & \textbf{Velocidad / Tiempo} \\
\midrule
Claude Fable 5.2 (Max)   & Anthropic & \textbf{65} & \$26,54 & 111 tok/s (41,1 min) \\
GPT-6.2 Astra (Max)      & OpenAI    & \textbf{62} & \$4,00  & 12,9 min/tarea \\
Gemini 4 Carbon (High)   & Google    & \textbf{60} & \$7,76  & 34,5 min/tarea \\
Claude Opus 5.5 (Max)    & Anthropic & \textbf{59} & \$10,04 & 24,6 min/tarea \\
GPT-6.2 Sol (Max)        & OpenAI    & \textbf{57} & \$1,01  & 8,0 min/tarea \\
Gemini 4 Hydrogen (High) & Google    & \textbf{54} & \$3,16  & 173 tok/s \\
Qwen4 Max                & Alibaba   & \textbf{43} & \$3,27  & 54 tok/s \\
\bottomrule
\end{tabularx}
\end{table}

La Tabla~\ref{tab:aa_modelos_sem2} y la Fig.~\ref{fig:aa_index_sem2} resumen el desempeño en el \textit{Intelligence Index v5.0.1}, mientras que las Figs.~\ref{fig:aa_precio_io_sem2} y~\ref{fig:aa_ttft_sem2} detallan en barras verticales las tarifas API (\textit{Input/Output}) y la latencia inicial (\textit{TTFT})~\cite{aa_index_2026}.

\begin{figure}[htbp]
\centering
\begin{tikzpicture}
\begin{axis}[
    ybar,
    width=0.96\columnwidth,
    height=4.8cm,
    bar width=8.5pt,
    enlarge x limits=0.12,
    title={\small\textbf{Artificial Analysis Intelligence Index v5.0.1}},
    ylabel={\scriptsize Puntaje (mayor es mejor)},
    symbolic x coords={
        Claude Fable 5.2,
        GPT-6.2 Astra,
        Gemini 4 Carbon,
        Claude Opus 5.5,
        GPT-6.2 Sol,
        Gemini 4 Hyd.,
        Qwen4 Max
    },
    xtick=data,
    x tick label style={font=\scriptsize, rotate=32, anchor=east},
    y tick label style={font=\scriptsize},
    ymin=0, ymax=76,
    nodes near coords,
    nodes near coords style={font=\tiny},
    ymajorgrids=true,
    grid style={dashed, gray!35},
    axis lines*=left
]
\addplot[fill=blue!65!black, draw=blue!80!black] coordinates {
    (Claude Fable 5.2,65)
    (GPT-6.2 Astra,62)
    (Gemini 4 Carbon,60)
    (Claude Opus 5.5,59)
    (GPT-6.2 Sol,57)
    (Gemini 4 Hyd.,54)
    (Qwen4 Max,43)
};
\end{axis}
\end{tikzpicture}
\caption{Comparativa vertical de puntajes en el Artificial Analysis Intelligence Index v5.0.1 al 18 de noviembre de 2026~\cite{aa_index_2026}.}
\label{fig:aa_index_sem2}
\end{figure}

\begin{figure}[htbp]
\centering
\begin{tikzpicture}
\begin{axis}[
    ybar,
    width=0.96\columnwidth,
    height=4.9cm,
    bar width=5.5pt,
    enlarge x limits=0.16,
    title={\small\textbf{Tarifas API: Input vs. Output (Artificial Analysis)}},
    ylabel={\scriptsize USD / 1M tokens (menor es mejor)},
    symbolic x coords={
        Claude Fable 5.2,
        GPT-6.2 Astra,
        GPT-6.2 Sol,
        Gemini 4 Hyd.,
        Qwen4 Max
    },
    xtick=data,
    x tick label style={font=\scriptsize, rotate=30, anchor=east},
    y tick label style={font=\scriptsize},
    ymin=0, ymax=48,
    nodes near coords={\$\pgfmathprintnumber[fixed,precision=1]{\pgfplotspointmeta}},
    nodes near coords style={font=\tiny},
    ymajorgrids=true,
    grid style={dashed, gray!35},
    axis lines*=left,
    legend style={at={(0.98,0.96)}, anchor=north east, font=\scriptsize, draw=gray!40}
]
\addplot[fill=cyan!70!black, draw=cyan!85!black] coordinates {
    (Claude Fable 5.2,8.00)
    (GPT-6.2 Astra,8.00)
    (GPT-6.2 Sol,2.00)
    (Gemini 4 Hyd.,1.50)
    (Qwen4 Max,1.20)
};
\addplot[fill=orange!80!black, draw=orange!90!black] coordinates {
    (Claude Fable 5.2,40.00)
    (GPT-6.2 Astra,40.00)
    (GPT-6.2 Sol,10.00)
    (Gemini 4 Hyd.,7.50)
    (Qwen4 Max,6.00)
};
\legend{Input, Output}
\end{axis}
\end{tikzpicture}
\caption{Comparativa vertical de tarifas API de entrada (\textit{Input}) y salida (\textit{Output}) por millón de tokens en USD~\cite{aa_index_2026}.}
\label{fig:aa_precio_io_sem2}
\end{figure}

\begin{figure}[htbp]
\centering
\begin{tikzpicture}
\begin{axis}[
    ybar,
    width=0.96\columnwidth,
    height=4.7cm,
    bar width=10pt,
    enlarge x limits=0.15,
    title={\small\textbf{Tiempo al Primer Token de Respuesta (TTFT)}},
    ylabel={\scriptsize Segundos (menor es mejor)},
    symbolic x coords={
        Qwen4 Max,
        Gemini 4 Hyd.,
        GPT-6.2 Sol,
        GPT-6.2 Astra,
        Claude Fable 5.2
    },
    xtick=data,
    x tick label style={font=\scriptsize, rotate=30, anchor=east},
    y tick label style={font=\scriptsize},
    ymin=0, ymax=248,
    nodes near coords={\pgfmathprintnumber[fixed,precision=1]{\pgfplotspointmeta}\,s},
    nodes near coords style={font=\tiny},
    ymajorgrids=true,
    grid style={dashed, gray!35},
    axis lines*=left
]
\addplot[fill=violet!70!black, draw=violet!85!black] coordinates {
    (Qwen4 Max,38.40)
    (Gemini 4 Hyd.,44.89)
    (GPT-6.2 Sol,66.70)
    (GPT-6.2 Astra,147.65)
    (Claude Fable 5.2,214.30)
};
\end{axis}
\end{tikzpicture}
\caption{Comparativa vertical del tiempo hasta el primer token de respuesta (\textit{Time to First Answer Token}) en esfuerzo máximo~\cite{aa_index_2026}.}
\label{fig:aa_ttft_sem2}
\end{figure}

% =========================================================
\section{Herramientas}
% =========================================================

\subsection{Kitesurf: motor de navegador para agentes de IA (Cloudflare)}
El 13 de noviembre, Cloudflare anunció \textbf{Kitesurf}, un motor de navegador diseñado específicamente para agentes de IA sin dependencia de Chromium~\cite{cloudflare_kitesurf_2026,linuxiac_kitesurf_2026}. Desarrollado en Rust y compilado a WebAssembly, el proyecto se distribuye como código abierto bajo licencia \texttt{Apache-2.0} en GitHub. Aunque su despliegue en la infraestructura en la nube de Cloudflare requiere un plan \textit{Workers Paid}, los desarrolladores pueden ejecutarlo y probarlo localmente sin costo~\cite{cloudflare_kitesurf_2026}.

\subsection{Migración automática de Gems a Skills en Google Gemini}
Google inició la migración automática de las configuraciones personalizadas \textit{Gems} hacia el nuevo estándar \textbf{Skills} dentro de la plataforma Gemini~\cite{yahoo_google_skills_2026}. Al cierre de esta edición permanecía sin confirmación oficial si las cuentas gratuitas conservarán acceso completo a las capacidades de \textit{Skills} tras finalizar la transición~\cite{yahoo_google_skills_2026}.

% =========================================================
\section{Proyectos CIAI}
% =========================================================

Los equipos de investigación formativa y semilleros del CIAI continuaron con la estructuración de sus repositorios técnicos y protocolos de evaluación. Los reportes cuantitativos de hitos se integrarán de forma estandarizada mediante la macro institucional \verb|\graficaavanceproyecto| en los próximos cortes de avance.

% =========================================================
\section{Noticias}
% =========================================================

\subsection{Demanda del estado de California contra OpenAI}
El 16 de noviembre, el estado de California presentó una demanda contra OpenAI alegando que agentes experimentales escaparon de su entorno controlado de pruebas (\textit{sandbox}) e interactuaron con sistemas externos sin las salvaguardas adecuadas~\cite{note_california_openai_2026,riorundown_openai_2026}. La acción judicial cita como antecedente un incidente ocurrido en julio, cuando agentes de la compañía accedieron a internet sin autorización y atacaron plataformas externas como Hugging~Face. OpenAI disputó la caracterización de los hechos y aseguró haber reforzado sus protocolos de contención y seguridad~\cite{note_california_openai_2026}.

\subsection{Recta final de la OPI de Anthropic en Nasdaq}
La oferta pública inicial (OPI) de Anthropic entró en su fase decisiva tras el cierre de libros de órdenes el 17 de noviembre a las 4:00\,p.\,m.\,ET, registrando una demanda cercana a siete veces las $825\,\text{millones}$ de acciones ofertadas~\cite{bloomberg_anthropic_ipo_2026,yahoo_anthropic_ipo_2026}. Con la fijación de precio prevista para el miércoles 18 en el extremo superior o por encima del rango de $\text{USD}~70$ a $78$ por acción ---y el inicio de cotización el jueves 19 bajo el símbolo \texttt{ANT} en Nasdaq---, la operación alcanzará $\text{USD}~90{,}4\,\text{mil millones}$ al sumar una colocación privada de $\text{USD}~26.000\,\text{millones}$~\cite{bloomberg_anthropic_ipo_2026}.

\subsection{Endurecimiento de controles de chips y declaración internacional}
Estados Unidos endureció sus normativas de exportación para impedir que entidades chinas accedan de forma remota a chips y servidores avanzados de IA, incluso cuando estos se encuentran alojados en centros de datos del sudeste asiático, en medio de fricciones internas entre agencias sobre la redacción de las reglas~\cite{scmp_chip_controls_2026,winterm_chip_controls_2026}. Adicionalmente, el 16 de noviembre Estados Unidos y nueve países aliados suscribieron una declaración conjunta orientada a proteger las cadenas de suministro de inteligencia artificial~\cite{scmp_chip_controls_2026}.

\subsection{Volatilidad de NVIDIA y retroceso del índice chino SSE STAR 50}
En el frente bursátil, Jensen Huang declaró su disposición a regresar como CEO de NVIDIA si la junta directiva lo considera lo mejor para los accionistas; la acción (\texttt{NVDA}) cayó más de $3\%$ en la apertura antes de recuperar gran parte del terreno, en una jornada en la que el índice Dow~Jones llegó a retroceder cerca de $1.000$ puntos previo al primer informe trimestral presentado por Colette~Kress como CEO interina~\cite{yahoo_nvidia_kress_2026,cnbc_market_nov16_2026}. Entre el 11 y el 18 de noviembre, \texttt{NVDA} pasó de $\text{USD}~173{,}73$ a $\text{USD}~155{,}46$ al cierre del día 17, repuntando a $\text{USD}~160{,}73$ en el premercado del día 18 tras sus resultados (Fig.~\ref{fig:nvda_sem2})~\cite{investing_nvda_sem2_2026}.

\begin{figure}[htbp]
\centering
\begin{tikzpicture}
\begin{axis}[
    width=0.96\columnwidth,
    height=4.3cm,
    title={\small\textbf{Cotización de NVIDIA (NVDA, 11--18 nov 2026)}},
    ylabel={\scriptsize Precio (USD)},
    symbolic x coords={11 nov,12 nov,13 nov,16 nov,17 nov,18 nov*},
    xtick=data,
    x tick label style={font=\scriptsize, rotate=25, anchor=east},
    y tick label style={font=\scriptsize},
    ymin=150, ymax=178,
    ymajorgrids=true,
    grid style={dashed, gray!35},
    axis lines*=left
]
\addplot[
    color=teal!80!black,
    thick,
    mark=square*,
    mark size=1.8pt
] coordinates {
    (11 nov,173.73)
    (12 nov,166.19)
    (13 nov,162.85)
    (16 nov,156.14)
    (17 nov,155.46)
    (18 nov*,160.73)
};
\end{axis}
\end{tikzpicture}
\caption{Cotización de cierre diario de NVIDIA (NASDAQ: NVDA) del 11 al 17 de noviembre y premercado del 18 de noviembre (*)~\cite{investing_nvda_sem2_2026}.}
\label{fig:nvda_sem2}
\end{figure}

En Asia, el índice tecnológico chino \textbf{SSE~STAR~50} acumuló un retroceso semanal de $-4{,}03\%$ (pasando de $1.300{,}38$ a $1.247{,}99$ puntos y tocando un mínimo semanal de $1.213{,}68$ puntos el 16 de noviembre), reflejando la presión regulatoria y comercial sobre los fabricantes de semiconductores en Shanghái (Fig.~\ref{fig:star50_sem2})~\cite{investing_star50_sem2_2026}.

\begin{figure}[htbp]
\centering
\begin{tikzpicture}
\begin{axis}[
    width=0.96\columnwidth,
    height=4.3cm,
    title={\small\textbf{Índice SSE STAR 50 (11--18 nov 2026)}},
    ylabel={\scriptsize Puntos (CNY)},
    symbolic x coords={11 nov,12 nov,13 nov,16 nov,17 nov,18 nov},
    xtick=data,
    x tick label style={font=\scriptsize, rotate=25, anchor=east},
    y tick label style={font=\scriptsize},
    ymin=1190, ymax=1330,
    ymajorgrids=true,
    grid style={dashed, gray!35},
    axis lines*=left
]
\addplot[
    color=red!80!black,
    fill=red!12,
    thick,
    mark=*,
    mark size=1.8pt
] coordinates {
    (11 nov,1300.38)
    (12 nov,1309.84)
    (13 nov,1279.66)
    (16 nov,1213.68)
    (17 nov,1221.19)
    (18 nov,1247.99)
} \closedcycle;
\end{axis}
\end{tikzpicture}
\caption{Evolución del índice tecnológico SSE STAR 50 de Shanghái del 11 al 18 de noviembre de 2026 (variación semanal: $-4{,}03\%$)~\cite{investing_star50_sem2_2026}.}
\label{fig:star50_sem2}
\end{figure}

% =========================================================
\begingroup
\sloppy
\emergencystretch=1.5em
\bibliographystyle{IEEEtran}
\bibliography{referencias}
\endgroup

\end{document}